%% kapa: exam=2 round=20 year=2009 period=3 subject=compensation_law subject_ko=감정평가 및 보상법규 seq=3 number=3 points=25
%% source: Q-Net 감정평가사 기출문제 (cst00309) artlSeq=5024180 file=extracted/1008/5024180/2009년 제20회 감정평가사 2차시험 기출문제/3교시.hwp
%% conversion: kordoc 4.24.51 -> tools/gen2.py (auto); review=proofread-sonnet; flags=none
\begin{longquestion}{【문제 3】}
「부동산 가격공시 및 감정평가에 관한 법률 시행령」제77조 제1항 [별표2] (감정평가업자의 설립인가의 취소와 업무의 정지에 관한 기준)는 재판규범성이 인정되는지의 여부를 설명하시오. (25점) \par
\medskip
\begin{framedtext}
\centering\textbf{【 회계관련 출제기준 적용 안내 】}\par
2010년 제21회 감정평가사 자격시험부터 회계처리 등과 관련된 시험문제는 한국채택국제회계기준(K-IFRS)만을 적용하여 출제하오니 수험준비에 참고하시기 바랍니다.
\end{framedtext}
\end{longquestion}
